% 02-marco-teorico.tex — Chapter 2: Marco Te\'{o}rico

\chapter{Marco Te\'{o}rico}
\label{ch:marco-teorico}

\section{Kernels de Markov y modelos de espacio de estados}

Un \textbf{kernel de Markov} $K: \mathcal{X} \times \mathcal{B}(\mathcal{X}) \to [0,1]$ es la probabilidad
condicional $K(x, A) = \Prob(X_t \in A \mid X_{t-1} = x)$. Generaliza las matrices de
transici\'{o}n a espacios de estado continuos \citep{chopin2020}.

Un \textbf{modelo de espacio de estados} (SSM) apila dos kernels:
\begin{align}
    X_t &\sim f(\cdot \mid X_{t-1}) \quad \text{(transici\'{o}n de estado)} \\
    Y_t &\sim g(\cdot \mid X_t) \quad \text{(emisi\'{o}n observacional)}
\end{align}
Las observaciones $\{Y_t\}$ \textbf{no} son Markov --- heredan dependencia a trav\'{e}s del
estado latente. Filtrar consiste en computar $P(X_t \mid Y_{1:t})$ integrando sobre toda
la trayectoria.

\section{Filtrado y suavizado}

El \textbf{filtrado} computa $P(X_t \mid Y_{1:t})$ usando solo datos hasta $t$.
El \textbf{suavizado} computa $P(X_t \mid Y_{1:T})$ usando toda la serie.
Para modelos lineales gaussianos, el \textbf{filtro de Kalman} da la soluci\'{o}n exacta:

\begin{align}
    \hat{x}_{t|t-1} &= A\hat{x}_{t-1}, \quad P_{t|t-1} = AP_{t-1}A^\top + Q \quad \text{(predicci\'{o}n)} \\
    \hat{x}_t &= \hat{x}_{t|t-1} + K_t(y_t - H\hat{x}_{t|t-1}) \quad \text{(actualizaci\'{o}n)}
\end{align}
donde $K_t$ es la ganancia de Kalman.

Para modelos no gaussianos (como SF-Harris), el filtro de Kalman no aplica.
La \textbf{formalizaci\'{o}n de Feynman-Kac}~\citep{chopin2020} proporciona el marco general:
el filtrado es una secuencia de integrales
$\int K(x_{t-1}, x_t) \cdot g(y_t \mid x_t) \cdot w(x_{t-1}) \, dx_{t-1}$,
y los filtros de part\'{\'i}culas las aproximan por muestreo de importancia secuencial.

\section{La cadena de Harris}

\begin{definicion}[Cadena de Harris]
Una cadena de Markov $\{X_t\}$ en $(\mathcal{X}, \mathcal{B})$ es \textbf{Harris-recurrente}
si existe un conjunto peque\~{n}o $C$ y una medida no trivial $\psi$ tales que
$\Prob(\tau_C < \infty \mid X_0 = x) = 1$ para todo $x \in \mathcal{X}$.
\end{definicion}

La cadena de Harris del modelo SF-Harris tiene kernel de transici\'{o}n:
\begin{equation}
    P(X_{t+1} \mid X_t) = \pstay \cdot \delta_{X_t} + (1 - \pstay) \cdot Q
\end{equation}
donde $\pstay = e^{-\alpha}$ es la probabilidad de estancia, $\delta_{X_t}$ es la
delta de Dirac en el estado actual, y $Q$ es la distribuci\'{o}n invariante.
La distribuci\'{o}n estacionaria es $Q$: el proceso converge independientemente del
valor inicial.

\section{La distribuci\'{o}n GIG}

La distribuci\'{o}n \textbf{Generalized Inverse Gaussian} (GIG)~\citep{barndorff1977}
tiene densidad:
\begin{equation}
    f(x \mid \lambda, \kappa, \eta) = \frac{(\kappa/\eta)^{\lambda/2}}{2 K_\lambda(\sqrt{\kappa\eta})} x^{\lambda-1} \exp\left(-\frac{\kappa x + \eta/x}{2}\right), \quad x > 0
\end{equation}
donde $K_\lambda$ es la funci\'{o}n de Bessel modificada de tercer tipo.
Casos especiales: $\lambda = 1$ da la distribuci\'{o}n hiperb\'{o}lica generalizada;
$\kappa \to 0$ da la gamma; $\eta \to 0$ da la inversa-gamma.

En el modelo SF-Harris original, $Q = \GIG(\lambda, \kappa, \eta)$.
Nuestros resultados muestran que la distribuci\'{o}n emp\'{\'i}rica supera a la GIG param\'{e}trica
(AAD 0.2\,pp vs 0.8\,pp).

\section{Muestreo de Gibbs para estimaci\'{o}n de par\'{a}metros}

El muestreador de Gibbs~\citep{chopin2020} produce muestras de la distribuci\'{o}n
posterior conjunta $P(\alpha, \mu, \sigma \mid Y_{1:T})$ muestreando iterativamente
de cada distribuci\'{o}n condicional completa:

\begin{enumerate}
    \item Muestrear $\alpha \mid \mu, \sigma, Y_{1:T}$ (metropolis dentro de Gibbs)
    \item Muestrear $\mu \mid \alpha, \sigma, Y_{1:T}$ (conjugado gaussiano)
    \item Muestrear $\sigma \mid \alpha, \mu, Y_{1:T}$ (conjugado inverso-gamma)
\end{enumerate}

Despu\'{e}s de $N_{\text{iter}}$ iteraciones y un per\'{\'i}odo de \emph{burn-in},
obtenemos muestras $\{(\alpha^{(i)}, \mu^{(i)}, \sigma^{(i)})\}_{i=1}^{N}$ que caracterizan
la incertidumbre param\'{e}trica.

\textbf{Observaci\'{o}n clave:} En nuestra implementaci\'{o}n, el muestreador de Gibbs
produce estimaci\'{o}n de par\'{a}metros (localizaci\'{o}n y escala de $Q$),
no suavizado del estado latente $\tau^*_{1:T}$.
La simulaci\'{o}n hacia adelante comienza desde la \'{u}ltima observaci\'{o}n ruidosa,
no desde un estado suavizado.